% Introduction (draft 2026-11-13). Numbers: tab1_controlled / tab5_reasoning / tab6_editing / tabC_multitarget,
% significance from results/metrics/sig_table1.json and sig_reason.json. Citation keys are placeholders:
%   ophrg 2607.15374, sam3 2511.16719, lisa 2308.00692, sa2va 2501.04001, attnseg 2610.18729, refcoco (Kazemzadeh
%   EMNLP 2014), refcocog_mao 1511.02283, paco 2301.01795, partimagenet 2112.00933, instructpart 2505.18291,
%   gres 2306.00968, gsva 2312.10103, affordjev 2611.07896, partgrounding_bottleneck 2609.13225.
\section{Introduction}
\label{sec:intro}

A request such as ``hold the mug by its handle'' or ``recolor the seat of the chair'' does not ask for every object
in the image that matches a phrase. It refers to one region: one part of one object. Answering it requires two
decisions that current language-guided segmentation systems tend to merge. The first is \emph{which} entity is
meant, the handle and not the mug, this chair and not the one behind it. The second is \emph{how much} of the image
that entity covers: the whole handle, not a sliver of it, and not the cup it is attached to. We call this task
\emph{part grounding}: given an image and a request, return the single mask of the part the request refers to.

Part grounding is where today's systems fail most visibly. Multimodal large language models (MLLMs) localize whole
objects well but return the parent object when a query names one of its parts~\cite{ophrg}; promptable segmenters
answer a noun phrase with every instance of the concept~\cite{sam3}; and token-based segmentation MLLMs decode a
mask from a single hidden state that carries no explicit notion of spatial extent~\cite{lisa,sa2va,attnseg}. In our
controlled comparison (Tab.~\ref{tab:controlled}), the strongest same-size pipeline, Qwen3-VL-4B boxes prompting
SAM\,3, reaches 33.0 gIoU on PACO parts against 76.4 on PACO objects.

The way these systems are usually evaluated hides part of the gap. Referring expressions were introduced to pick
out one object~\cite{refcoco,refcocog_mao}, and part and affordance benchmarks pair each request with one target
region~\cite{paco,partimagenet,instructpart}. Yet recent protocols score the union of every instance a method
returns~\cite{ophrg,gres}. Under a union rule, covering the whole object or several guesses still earns overlap:
with all instances merged, SAM\,3 reaches 53.9 gIoU on PartImageNet~\cite{ophrg}, while its single top-scoring
mask at its released threshold reaches 24.3 (Tab.~\ref{tab:controlled}). Downstream uses do not have this slack.
An editor applies the mask it is given, and a robot grasps the part it is shown~\cite{affordjev,partgrounding_bottleneck}; a wrong or additional region edits or touches the
wrong pixels. We therefore evaluate every method under one rule: \emph{at most one mask per query}, the answer the
request asks for. Multi-target expressions~\cite{gres,gsva} remain a natural extension, which we evaluate
separately.

Our approach starts from an observation about where the difficulty lies. A promptable segmenter already proposes
a good mask for most parts, but does not rank it first: on PARTONOMY, the best of the candidates SAM\,3 produces for
a query reaches 71.0 gIoU, while SAM\,3's own top-ranked instance reaches 42.0 (Tab.~\ref{tab:coverage}). What is
missing is the decision among them, which requires reading the request in the context of the whole image. We give
that decision to an MLLM. We propose \ours{}, which \emph{asks} the MLLM to choose a mask instead of \emph{drawing}
one. For each query, \ours{} builds a pool of candidate masks from a frozen SAM\,3: automatic masks from a point
grid, which are query-independent and cover objects and parts at every scale, and the top instances SAM\,3 returns
for the query's own phrases. Every candidate enters the MLLM as a \emph{mask token}, an embedding pooled from the
image tokens it covers together with its geometry and its provenance. Each mask token attends only to the image
tokens of its own region, so it represents that region rather than the whole image. A selection token at the end
of the query reads the pool and scores every candidate; the answer is the highest-scoring mask, or no mask when no
candidate passes a threshold chosen on validation (needed only where a request can have no target).

This design separates the two decisions that part grounding requires. The extent of each answer comes from a
segmentation model built for boundaries, and the choice of answer comes from a language model built for reading
requests in context. The MLLM never generates mask pixels, so its reasoning is spent on the choice alone. \ours{}
adds 64.2M trainable parameters (1.2\% of the system): LoRA adapters on Qwen3-VL-4B, the mask-token embedding and
the selection heads, trained on PACO, Pascal-Part and gRefCOCO with SAM\,3 frozen.

Under the one-mask protocol, with Qwen3-VL-4B for every MLLM-based method and test splits throughout, \ours{}
improves part grounding on five of six benchmarks over the strongest same-size baseline, with paired image-level
bootstrap tests at $p<10^{-4}$: $+14.5$ gIoU on PACO parts, $+13.9$ on GSEval parts, $+7.1$ on PascalPart-116,
$+5.6$ on PartImageNet and $+2.3$ on PARTONOMY, and it is on par on InstructPart part names
(Tab.~\ref{tab:controlled}). The same model improves referring segmentation on gRefCOCO by $+8.0$ / $+7.9$ gIoU and
rejects expressions with no target more reliably ($+4.9$ / $+8.1$ N-acc). On reasoning queries, where a part is
described by the task it serves rather than named, \ours{} with query rewriting by the same backbone exceeds
Box + SAM\,3 by $+11.2$ / $+8.5$ gIoU on InstructPart and is on par on the object-level ReasonSeg
(Tab.~\ref{tab:reasoning}). As a mask source for region-grounded editing, \ours{} returns a mask for every request
and is on par with same-size baselines in edit quality, with the best mask IoU on ImgEdit (Tab.~\ref{tab:editing}).
When several masks per expression are allowed, the same scores select a set and lead gRefCOCO's multi-target protocol
(Tab.~\ref{tab:multitarget}).

In summary, our contributions are:
\begin{itemize}
\item We formulate \emph{part grounding} as selecting one region among candidate masks, and evaluate it under a
  uniform one-mask-per-query protocol that every method, including ours, follows; we report the multi-mask setting
  separately.
\item We propose \ours{}, which represents each candidate mask as a region-restricted token inside an MLLM and
  answers a request by selection, separating the choice of region (the MLLM) from its extent (a frozen segmenter)
  with 64.2M trainable parameters.
\item In a controlled comparison on test splits with a same-size backbone, \ours{} improves part grounding on five of
  six benchmarks, referring segmentation and rejection on gRefCOCO, and part-level reasoning, and supports
  region-grounded editing with a mask for every request; ablations isolate query-conditioned candidates, mask
  tokens and region-restricted attention.
\end{itemize}
